\section{From masked energies to the maximal transform}\label{sec:maximal}

The estimate of Section~\ref{sec:jumps} allows different truncation
endpoints at every output position. We first use this freedom on a finite
mesh, choosing the endpoints that maximize each output. We then pass to
functions and compare the resulting smooth kernels with hard truncations.
The comparison will be pointwise and uniform in both endpoints.

\subsection{Odd windows and selected endpoints}

For $s>0$ and $0<\varepsilon<R$, define
\[
 k_s(t)=\frac{t}{\sqrt{s}}g_s(t),\qquad
 \psi=k_1*k_1*k_1,\qquad
 I_{\varepsilon,R}(t)=\int_{\varepsilon^2}^{R^2}
       s^{-1/2}\psi(t/\sqrt{s})\,\frac{\dd s}{s}.
\]
Here $*$ is convolution on $\R$. The function $\psi$ is Schwartz, and
$I_{\varepsilon,R}$ is smooth and integrable for fixed positive endpoints.

\begin{proposition}[Selected smooth truncations]\label{prop:selected-kernel}
Let $h>0$, let $N\ge0$ be an integer, set $u_i=ih$, and let
$A_v(i,j)$, $v=0,1,2$, be complex
arrays indexed by $-N\le i,j\le N$. Assign an arbitrary endpoint pair
$h\le\varepsilon_{ij}<\rho_{ij}<\infty$ to each $(i,j)$. Then
\begin{equation}\label{eq:selected-kernel}
 \left|h^3\sum_{i,j,l}A_0(i,j)A_1(j,l)A_2(l,i)
       I_{\varepsilon_{ij},\rho_{ij}}(u_i+u_j+u_l)\right|
 \le C\prod_{v=0}^2\left(\sum_{i,j}h^2|A_v(i,j)|^3\right)^{1/3}.
\end{equation}
The constant is independent of the arrays, mesh and selected endpoints.
\end{proposition}

\begin{proof}
We apply the mixed trace inequality with Gaussian weights of larger
variance than the windows $k_s$, so that the diagonal insertions are
bounded. We then compare each quadratic cost with a variance choice
covered by Proposition~\ref{prop:masked-dissipation}.

Use the masked array
$A_0^s(i,j)=A_0(i,j)\mathbf1_{\{\varepsilon_{ij}^2<s<\rho_{ij}^2\}}$
and the scale interval $[s_-,s_+]$ from
Proposition~\ref{prop:masked-dissipation}. All weighted matrices below use
this same mask. Let superscript $\mathrm b$ denote the base variances
$(\lambda s,\lambda s,\lambda s)$, where $\lambda=11/10$, and put
\[
 D_v(i,i)=\frac{k_s(u_i-p_v)}{g_{\lambda s}(u_i-p_v)}.
\]
These matrices are real diagonal. With $z=(u_i-p_v)/\sqrt{s}$,
\[
 |D_v(i,i)|=\sqrt\lambda\,|z|
       e^{-(\lambda-1)z^2/(2\lambda)}
 \le\frac{\lambda}{\sqrt{e(\lambda-1)}}.
\]
Thus their operator norms have an absolute bound.

For each vertex $v$, let $\boldsymbol\alpha(v)$ equal $1$ at $v$ and
$\lambda$ at the other two vertices. In the star centered at $v$, changing
from the base variances to $\boldsymbol\alpha(v)s$ changes only its row
weight. In particular,
\[
 R_v^{\boldsymbol\alpha(v)}
 =\diag_i\left(\sqrt{\frac{g_s(u_i-p_v)}
                             {g_{\lambda s}(u_i-p_v)}}\right)R_v^{\mathrm b},
 \qquad
 T_v^{\boldsymbol\alpha(v)}\le\sqrt\lambda\,T_v^{\mathrm b}.
\]
The last inequality uses $g_s\le\sqrt\lambda\,g_{\lambda s}$.
The score $L_v^{\boldsymbol\alpha(v)}$ equals
$\diag_i((u_i-p_v)/s)$, so the row densities also give the exact identity
\begin{equation}\label{eq:maximal-window-score}
 (R_v^{\mathrm b})^*D_vR_v^{\mathrm b}
 =\sqrt{s}\,(R_v^{\boldsymbol\alpha(v)})^*
           L_v^{\boldsymbol\alpha(v)}R_v^{\boldsymbol\alpha(v)}.
\end{equation}
Indeed, both inserted row densities are $h k_s(u_i-p_v)$, and all column
weights agree. The insertion on the left is therefore supported on
$T_v^{\boldsymbol\alpha(v)}$. Lemma~\ref{lem:order} and quadratic
homogeneity give
\[
 b_{T_v^{\mathrm b}}\bigl((R_v^{\mathrm b})^*D_vR_v^{\mathrm b}\bigr)
 \le\lambda^{1/4}s V_v^{\boldsymbol\alpha(v)}.
\]
Apply Lemma~\ref{lem:mixed-trace} to the base matrices. Integrate its three
nonnegative terms separately, using Proposition~\ref{prop:masked-dissipation}
with the corresponding $\boldsymbol\alpha(v)$. The factor $s$ cancels
$\dd s/s$, and we obtain
\begin{equation}\label{eq:masked-window-integral}
 \int_{s_-}^{s_+}\int_{\Pi_0}
 \left|\tr(D_0W_0^{\mathrm b}D_1W_1^{\mathrm b}D_2W_2^{\mathrm b})\right|
       \,\dd\pi_0\,\frac{\dd s}{s}
 \le C\mathcal S(A).
\end{equation}

We identify the signed integral in this estimate. Expanding the trace
leaves one signed Gaussian weight at each vertex:
\[
 h^3\sum_{i,j,l}A_0^s(i,j)A_1(j,l)A_2(l,i)
       k_s(u_i-p_0)k_s(u_j-p_1)k_s(u_l-p_2).
\]
On $\Pi_0$, substitute $r=u_i-p_0$ and $q=u_j-p_1$.
The Jacobian is one and the third argument becomes
$u_i+u_j+u_l-r-q$. Consequently the center integral is
$(k_s*k_s*k_s)(u_i+u_j+u_l)$. Since
$k_s(t)=s^{-1/2}k_1(t/\sqrt{s})$, this convolution equals
\[
 s^{-1/2}\psi\bigl((u_i+u_j+u_l)/\sqrt{s}\bigr).
\]
The sums are finite and the Gaussian products are absolutely integrable
on the center plane and the compact positive scale interval. Thus all
these integrations can be performed term by term. The mask restricts the
scale integral of the $(i,j,l)$ term to
$(\varepsilon_{ij}^2,\rho_{ij}^2)$. Taking the absolute value after
integration in \eqref{eq:masked-window-integral} proves
\eqref{eq:selected-kernel} with its right side initially replaced by
$C\mathcal S(A)$.

Put $n_v=(\sum h^2|A_v|^3)^{1/3}$. If any $n_v$ is zero the signed sum
vanishes. Otherwise apply the last estimate to $A_v/n_v$ for all three
vertices, keeping the selected endpoints fixed. The normalized value of
$\mathcal S$ is $3$. Scaling the signed trilinear expression back proves
the proposition.
\end{proof}

\subsection{Finite maxima and the continuum limit}

For functions $F,G$ on $\R^2$, write
\[
 B^I_{\varepsilon,R}(F,G)(x,y)
   =\int_\R F(x+t,y)G(x,y+t)I_{\varepsilon,R}(t)\,\dd t.
\]
For compactly supported smooth inputs this integral is defined at every
point. We now place the supremum over its endpoints inside the spatial
norm.

\begin{proposition}\label{prop:smooth-kernel-maximal}
For all complex $F,G\in C_c^\infty(\R^2)$,
\[
 \left\|\sup_{0<\varepsilon<R}
          |B^I_{\varepsilon,R}(F,G)|\right\|_{L^{3/2}(\R^2)}
 \le C\|F\|_{L^3(\R^2)}\|G\|_{L^3(\R^2)}.
\]
The supremum is measurable and the constant is absolute.
\end{proposition}

\begin{proof}
First fix a nonempty finite family of endpoint pairs
$(\varepsilon_k,R_k)$, $1\le k\le m$, and take
$h\le\min_k\varepsilon_k$. For finite arrays $A_1,A_2$, define
\[
 U_k(i,j)=h\sum_l A_1(j,l)A_2(l,i)
                  I_{\varepsilon_k,R_k}(u_i+u_j+u_l),
 \qquad M(i,j)=\max_{1\le k\le m}|U_k(i,j)|.
\]
At each output pair choose the least index attaining this maximum, and
write $U(i,j)$ for its output. These choices provide admissible
$\varepsilon_{ij},\rho_{ij}$ in Proposition~\ref{prop:selected-kernel}.
They are fixed while the testing array $A_0$ varies.

Set $H=\sum_{i,j}h^2 M(i,j)^{3/2}$. If $H=0$ the desired bound is
immediate. Otherwise choose
\[
 A_0(i,j)=
 \begin{cases}
 H^{-1/3}M(i,j)^{1/2}\,\overline{U(i,j)}/|U(i,j)|,
       &U(i,j)\ne0,\\
 0,&U(i,j)=0.
 \end{cases}
\]
Then $\sum h^2|A_0|^3=1$ and $h^2\sum A_0U=H^{2/3}$.
Proposition~\ref{prop:selected-kernel} therefore gives
\begin{equation}\label{eq:discrete-finite-maximum}
 \left(\sum_{i,j}h^2M(i,j)^{3/2}\right)^{2/3}
 \le C\prod_{v=1}^2\left(\sum_{i,j}h^2|A_v(i,j)|^3\right)^{1/3}.
\end{equation}
This estimate has no dependence on the number of endpoint pairs. The
complex phase of $A_0$ accounts for the absolute value at each output.

To pass to functions, define the compactly supported smooth inputs
\[
 G_1(b,c)=F(b+c,-b),\qquad G_2(c,a)=G(-a,c+a).
\]
Both pair maps have determinant one, so
$\|G_1\|_3=\|F\|_3$ and $\|G_2\|_3=\|G\|_3$.
Choose $B>0$ so that both supports lie in the interior of $[-B,B]^2$.
Take $h=B/N$, $-N\le i,j,l\le N$, and sample
\[
 A_1(j,l)=G_1(jh,lh),\qquad A_2(l,i)=G_2(lh,ih).
\]
The corresponding continuum outputs are
\[
 \mathcal U_k(a,b)=\int_\R G_1(b,c)G_2(c,a)
                         I_{\varepsilon_k,R_k}(a+b+c)\,\dd c.
\]
They vanish outside $[-B,B]^2$: $G_2(c,a)$ restricts $a$ and $G_1(b,c)$
restricts $b$. A nonzero integrand also has $c\in[-B,B]$.
For our fixed finite endpoint family the integrands and their $c$
derivatives are uniformly bounded on this compact cube. Riemann sums
in $c$ consequently converge uniformly for $(a,b)\in[-B,B]^2$ and
$1\le k\le m$. In particular,
\[
 \max_{i,j}\max_{1\le k\le m}
       |U_k(i,j)-\mathcal U_k(ih,jh)|\longrightarrow0.
\]
The finite maximum of magnitudes preserves this uniform convergence.
Since $\max_k|\mathcal U_k|$ is continuous and compactly supported,
the two-dimensional Riemann sums on the left of
\eqref{eq:discrete-finite-maximum} converge to its $L^{3/2}$ norm.
The cube roots of the input sums converge to $\|G_1\|_3$ and
$\|G_2\|_3$.
The condition on $h$ holds for all sufficiently large $N$.
We obtain
\[
 \|\max_k|\mathcal U_k|\|_{3/2}\le C\|G_1\|_3\|G_2\|_3.
\]
No convergence or regularity of the maximizing indices is needed: we
passed to the limit in the finite maximum itself.

With $x=-a$, $y=-b$, and $t=a+b+c$, the last integral is
$B^I_{\varepsilon_k,R_k}(F,G)(x,y)$. The output coordinate change
preserves Lebesgue measure, giving the asserted estimate for every
finite endpoint family.

Finally enumerate all rational pairs $0<\varepsilon<R$ and take
increasing finite subfamilies. At each spatial point the output is
continuous in the two positive endpoints, by its defining integrals
on compact scale and spatial intervals. Thus the supremum over these
rational pairs equals the supremum over all real pairs. It is measurable
as a countable supremum of continuous spatial functions. Monotone
convergence of the $3/2$ powers of the finite maxima proves the result.
\end{proof}

\subsection{A pointwise comparison with hard truncations}

The smooth maximal estimate will imply the hard one because each
endpoint error has a common pointwise majorant. A uniform bound on
its kernel's $L^1$ norm alone would not provide this comparison.

\begin{lemma}\label{lem:maximal-hard-kernel}
There are an absolute constant $C$ and functions $E_r$, $r>0$, such that,
for $t\ne0$, apart from endpoint values,
\begin{equation}\label{eq:maximal-hard-decomposition}
 I_{\varepsilon,R}(t)
  =\frac{c_0}{t}\mathbf1_{\{\varepsilon<|t|<R\}}
       +E_\varepsilon(t)-E_R(t),
 \qquad c_0=-\frac{2}{3\sqrt{6\pi}},
\end{equation}
and
\begin{equation}\label{eq:maximal-error-envelope}
 |E_r(t)|\le\frac{C}{r}\left(1+\frac{|t|}{r}\right)^{-2}
 \qquad(t\ne0).
\end{equation}
\end{lemma}

\begin{proof}
Since $k_1=-g_1'$ and Gaussian variances add under convolution,
\[
 \psi(t)=-g_3'''(t)
        =\left(\frac{t^3}{27}-\frac t3\right)g_3(t).
\]
The convolution identity follows by completing the square;
differentiation under convolution is justified by Schwartz decay.
For $v\ge0$, put
\[
 P(v)=2\int_0^v\psi(z)\,\dd z
     =2g_3''(0)-2g_3''(v),\qquad
 c_0=P(\infty)=2g_3''(0)=-\frac2{3\sqrt{6\pi}}.
\]
The substitution $z=t/\sqrt{s}$ for $t>0$, followed by oddness for
$t<0$, yields
\[
 I_{\varepsilon,R}(t)
   =\frac{P(|t|/\varepsilon)-P(|t|/R)}{t}.
\]
Define $Q(v)=P(v)-c_0\mathbf1_{\{v>1\}}$ for $v\ge0$.
Smooth oddness of $\psi$ gives $P(v)=O(v^2)$ at zero. For $v>1$,
\[
 Q(v)=-2g_3''(v)
      =-2\left(\frac{v^2}{9}-\frac13\right)g_3(v).
\]
Thus $Q$ has Gaussian decay at infinity and bounded one-sided values
at its sole jump, $v=1$. Enlarging an absolute constant on a compact
interval, we have
\[
 \frac{|Q(v)|}{v}\le C(1+v)^{-2}\qquad(v>0).
\]
Set $E_r(t)=Q(|t|/r)/t$ for $t\ne0$ and set $E_r(0)=0$.
Substituting the definition of $Q$ proves
\eqref{eq:maximal-hard-decomposition} and
\eqref{eq:maximal-error-envelope}. Values at the four hard endpoints
have no effect on any integral.
\end{proof}

Let $M_x$ and $M_y$ denote the continuous centered interval maximal
operators in the first and second variables, respectively. The next
estimate follows from their one-dimensional bounds in
Lemma~\ref{lem:line-maximal}.

\begin{lemma}\label{lem:pointwise-endpoint-error}
For complex $F,G\in C_c^\infty(\R^2)$ and every $r>0$,
\[
 \left|\int_\R F(x+t,y)G(x,y+t)E_r(t)\,\dd t\right|
 \le C\bigl[M_x(|F|^2)(x,y)M_y(|G|^2)(x,y)\bigr]^{1/2}.
\]
The right side has $L^{3/2}(\R^2)$ norm at most $C\|F\|_3\|G\|_3$.
\end{lemma}

\begin{proof}
Put $q_r(t)=r^{-1}(1+|t|/r)^{-2}$. For nonnegative locally integrable
$f$ on $\R$, the part of $\int q_r(t)f(x+t)\,\dd t$ with $|t|\le r$
is at most $2M_{\R}f(x)$. On the shell
$2^kr<|t|\le2^{k+1}r$, $k\ge0$, its contribution is at most
\[
 r^{-1}2^{-2k}\int_{|t|\le2^{k+1}r}f(x+t)\,\dd t
 \le 4\,2^{-k}M_{\R}f(x).
\]
Summing the geometric series proves a bound independent of $r$.
By \eqref{eq:maximal-error-envelope} and weighted Cauchy--Schwarz,
\begin{align*}
 \left|\int F(x+t,y)G(x,y+t)E_r(t)\,\dd t\right|
 &\le C\left(\int |F(x+t,y)|^2q_r(t)\,\dd t\right)^{1/2}\\
 &\hspace{12mm}\cdot
       \left(\int |G(x,y+t)|^2q_r(t)\,\dd t\right)^{1/2}.
\end{align*}
Apply the preceding one-dimensional bound on the respective coordinate
lines. This gives the pointwise assertion.

H\"older's inequality and Lemma~\ref{lem:line-maximal}, applied on slices
and then integrated in the other coordinate, give
\begin{align*}
 \|[M_x(|F|^2)M_y(|G|^2)]^{1/2}\|_{3/2}
 &\le\|M_x(|F|^2)\|_{3/2}^{1/2}
       \|M_y(|G|^2)\|_{3/2}^{1/2}\\
 &\le C\||F|^2\|_{3/2}^{1/2}\||G|^2\|_{3/2}^{1/2}
  =C\|F\|_3\|G\|_3.
\end{align*}
\end{proof}

\begin{proposition}[Compact smooth inputs]\label{prop:compact-maximal}
For all complex $F,G\in C_c^\infty(\R^2)$,
\[
 \|B_*(F,G)\|_{L^{3/2}(\R^2)}
 \le C\|F\|_{L^3(\R^2)}\|G\|_{L^3(\R^2)},
\]
with an absolute constant. The supremum defining $B_*$ is over both
finite hard truncation endpoints, inside the spatial norm.
\end{proposition}

\begin{proof}
Lemma~\ref{lem:maximal-hard-kernel} and
Lemma~\ref{lem:pointwise-endpoint-error} give, at each spatial point,
\[
 |c_0|\sup_{0<\varepsilon<R}|B_{\varepsilon,R}(F,G)|
 \le \sup_{0<\varepsilon<R}|B^I_{\varepsilon,R}(F,G)|
       +2C[M_x(|F|^2)M_y(|G|^2)]^{1/2}.
\]
Both endpoint errors have the same majorant, independently of their
radii. The hard outputs are continuous in both positive endpoints,
so their supremum also equals the measurable supremum over rational
pairs. Take $L^{3/2}$ norms, apply
Proposition~\ref{prop:smooth-kernel-maximal} and
Lemma~\ref{lem:pointwise-endpoint-error}, and divide by $|c_0|>0$.
\end{proof}
